Experiments are conducted on the DataSense CIC IIoT 2025 dataset using 10-fold stratified cross-validation \cite{firouzi2025datasense}. Within each fold, the scaler is fitted on the training partition, the VICReg encoder is trained using only that partition, feature selection is performed on the corresponding training representation, and the downstream classifier is evaluated on the untouched test partition. \rev{At its natural operating point, the method retains on average 59 of the 71 numeric traffic features and achieves macro-F1 scores of 0.936 for binary detection and 0.908 for eight-class attack-category classification with Random Forest. A paired equivalence test confirms that this performance lies within $\pm$0.01 macro-F1 of the all-feature configuration, and the method significantly outperforms MCFS and the supervised DataSense-17 subset. The evaluation is complemented by three independent master seeds, device- and attack-execution-grouped folds, controlled component comparisons, and two additional public IoT/IIoT datasets, N-BaIoT and WUSTL-IIoT-2021.}

Overall, \method{} is intended for settings in which preserving unlabeled behavioral structure and supporting subsequent interpretation are more important than minimizing feature cardinality alone. Its cluster profiles\rev{, which are stable across folds and seeds and are reproduced by the proxy on unseen data,} provide behavioral evidence related to traffic characteristics such as TCP-flag activity, time-to-live statistics, and packet-size distributions, while remaining independent of attack labels during representation learning and feature selection.
